第 2 部の主題は、検査の精度を統合するメタ解析です。問いを立て、文献を探して選び、データを取り出し、バイアスのリスクを評価し、統合して報告するという工程は、第 1 部の介入のメタ解析と同じです。違うのは統計の部分で、検査の研究は感度と特異度という互いに関係した 2 つの結果を持つため、介入のときのように 1 つの効果の大きさを平均するわけにはいきません。

第 9 章と第 10 章では、精度を表す言葉と、検査のメタ解析で二変量モデルを使う理由を説明します。第 10 章では、感度と特異度を logit の目盛りに直す考え方が出てきます。第 11 章と第 12 章では、架空のプロカルシトニン（PCT）の研究を例に、抄録の選別、2×2 表の取り出し、バイアスのリスクの評価を練習します。第 13 章から第 15 章で、R による解析と結果の報告に進みます。第 1 部の第 1 章から第 5 章を読んでいることを前提にしますが、第 9 章と第 10 章だけを読んでも、検査のメタ解析の論文を読むための知識は一通り身につきます。第 2 部に入る前に、対数とオッズの\kkref{sec:05-1}、固定効果とランダム効果の\kkref{sec:05-4}、I² の\kkref{sec:06-1}、予測区間の\kkref{sec:06-3}、PRISMA のフロー図の\kkref{sec:03-5}を読み直してください。発展の印が付いた部分は、初読では飛ばしても後の章が読めます。
